\section{Datenmodellierung}
alle relevanten Informationen eines Systems beschreiben – mit einer Syntax, die einfach,
aber formal/exakt genug ist. Schnittstelle zwischen Realität und konzeptionellem Modell.
\subsection{Grundbegriffe}
\begin{itemize}
  \item \textbf{Objekt (entity)} ein konkretes Exemplar (z. B. Lehrer „Fröhlich“).
  \item \textbf{Objekttyp (entity type)} die Klasse gleichartiger Objekte (z. B. Lehrer).
  \item \textbf{Eigenschaft (attribute)} Merkmal aller Objekte; 
  \item  \textbf{Eigenschaftswert} konkrete Ausprägung
  \item  \textbf{Beziehung (relationship)} Zusammenhang zwischen zwei Objekttypen.
  \item \textbf{Beziehungstyp} alle gleichartigen Beziehungen.
\end{itemize}
\subsection{Stufen der Datenmodellierung}
\begin{itemize}
  \item \textbf{Konzeptionell} Fachliche Beschreibung, Identifikation von Entitäten,Attributen,Beziehungen – meist als ER-Diagramm – aus den Bestehnden Daten. Die Phase ist rein Konzeptionell deshalb werden noch keine Datentypen genutzt.
  Es wird drauft geachtet wo fallen daten an und Welche daten sind relevanten
  \item \textbf{Logisch} Anhand des ER-Diagramms Tabellen erstellen, Festlegung von Primär-/Fremdschlüsseln, Normalisierung. Erstellung eines relationalen Schema. 
  \item \textbf{Physikalisch}  Umsetzung der Tabbelen in einem konkreten DBMS: Datentypen, Indizes, Speicherstrukturen. Daruas folgt ein implementiertes Schema. Also auch SQL "Code"
\end{itemize} 